% !PRINTED 243
\weberpage{243}
In order to gain a foundation for this theory, we consider an integral function of a variable quantity $t$

\weq{1}{f(t) = a_0 t^n + a_1 t^{n-1} + \cdots + a_{n-1} t + a_n,}

whose coefficients are to be integers. We choose in addition a prime $p$ as modulus, and assume that $a_0$ is not divisible by $p$.

Two integral functions of $t$ in which corresponding coefficients are congruent with respect to the modulus $p$ shall themselves be called \gesperrt{congruent with respect to the modulus $p$}.

The function $f(t)$ is called \gesperrt{reducible with respect to the modulus $p$} if there exist two integral functions $\varphi(t)$, $\psi(t)$ with integral coefficients, in each of which at least one term depending on $t$ has a coefficient not divisible by $p$, such that

\weq{2}{f(t) \equiv \varphi(t)\,\psi(t)\ (\mathrm{mod}\ p)}

holds. In $\varphi(t)$ and $\psi(t)$ one may omit all terms whose coefficients are divisible by $p$, and then the degree of $\varphi(t)\,\psi(t)$ must agree with the degree of $f(t)$. Hence the degrees of $\varphi(t)$ and $\psi(t)$ must be lower than $n$. If there are no such functions $\varphi(t)$, $\psi(t)$, then $f(t)$ is called \gesperrt{irreducible with respect to the modulus $p$}.

A function irreducible with respect to the modulus $p$ is certainly always absolutely irreducible in the field of rational numbers. But the converse need not hold.

Thus, for example, the function of degree $2$, $t^2 - R$, is reducible with respect to $p$ if $R$ is a quadratic residue of $p$; for if $a^2 \equiv R\ (\mathrm{mod}\ p)$, then

\[
t^2 - R \equiv (t-a)(t+a)\ (\mathrm{mod}\ p).
\]

On the other hand $t^2 - N$, if $N$ is a non-residue of $p$, is irreducible, because otherwise $t^2 - N$ would have to become divisible by $p$ for a rational $t$.

Another example is offered by the functions

\[
t^2 + t + 1, \quad t^3 + t + 1,
\]

which are irreducible for the modulus $2$. For were they reducible, one of the factors would have to be linear, and there would have to be an integer which, substituted for $t$, makes these functions even numbers. But that is evidently impossible, since both functions are odd numbers for $t = 0$ and $t = 1$.
